%============================================================ 4. 제안 방법
\section{제안 방법}
\label{sec:method}

\subsection{판단 지연이 다른 2계층 비동기 구조}
\label{sec:arch}

제안 체계는 방어를 위한 판단을 전략 계층과 기동 계층, 두 계층으로 분리한다(Fig.~\ref{fig:architecture}). 
\textbf{두 계층은 같은 주기로 갱신되지만 응답 지연이 크게 다르며}, 느린 쪽을 비동기로 분리한 것이 구조의 핵심이다.

%-------------------------------------------------- 그림 2
\begin{figure*}[t]
\centering
\includegraphics[width=\textwidth]{fig2_architecture.pdf}
\caption{2계층 비동기 구조와 갱신 주기. 두 판단 계층은 \emph{같은} 결정 주기
         (25\,step $=$ 25초)로 갱신되지만 응답 지연이 다르다 --- 전략 계층은 로컬 14B
         모델에서 평균 16.6초로 주기의 3분의 2를 소모하는 반면, 기동 계층의 순전파 시간은
         그에 비해 무시할 수 있다. 그래서 전략 계층 호출만 비동기로 분리한다. 저수준
         제어(PD 조타)는 그보다 25배 빠른 1초 주기로 돈다.}
\label{fig:architecture}
\end{figure*}

전략 계층은 \emph{어느 방어정이 어느 적 무리를 맡는가}를 정하고, 기동 계층은 \emph{배정된
무리를 막기 위해 해면 어디에 그물벽을 놓는가}를 정한다. 시뮬레이션 시간 간격이
$\Delta t = 1\unit{s}$ 이고 결정 주기가 25\,step 이므로, 한 결정은 실시간 25초에 해당한다.
\textbf{두 계층은 모두 이 주기로 갱신된다.} 상위 계층의 재계획 주기를 늘리면 API 호출은
줄지만 배정 계획이 낡아 비교가 왜곡되므로(\S\ref{sec:experiment}), 본 연구는 두 계층을
같은 주기로 운용한다.

계층을 가르는 것은 주기가 아니라 \emph{지연}이다. 기동 계층의 순전파는 결정당 1회로 끝나는
반면, 전략 계층의 언어모델 호출은 \textbf{로컬 14B 모델에서} 평균 16.6초가 걸려 한 결정
주기의 3분의 2를 소모한다(Table~\ref{tab:planquality}). 동기로 기다리면 남는 여유가 8초
남짓이고, 실측에서 그 호출의 4.2\,\%가 주기 전체를 넘겼다(\S\ref{sec:res_latency}). 그래서
전략 계층 호출만 비동기로 분리한다. 호출이 진행되는 동안 기동 계층은 직전 배정을 그대로
사용하여 결정을 계속하고, 응답이 도착하면 다음 결정부터 새 배정이 반영된다. 다만 이 지연은
지휘관 모델에 따라 크게 다르며, 클라우드 모델과 소형 로컬 모델에서는 주기의 10\,\% 내외에
그친다(\S\ref{sec:res_latency}).

한편 저수준 제어(PD 조타)는 매 tick, 즉 1초마다 동작한다. 따라서 체계의 갱신 주기는
25초(두 판단 계층)와 1초(제어)의 \emph{두 가지}이며, 같은 결정 주기 안에서 두 판단 계층의
응답 지연이 다르다.

\subsection{기동 계층 (1): 기하 휴리스틱 기준 방법}
\label{sec:heuristic}

기동 계층을 학습으로 대체하기 전에, 같은 문제를 기하만으로 푸는 기준 방법을 먼저 정의한다.
이 방법은 두 가지 역할을 한다. 첫째 \S\ref{sec:results}의 비교 기준이고, 둘째 학습 초기의
행동 목표(\S\ref{sec:training})이다. 이 절의 두 알고리즘은 언어모델을 전혀 쓰지 않는다.

\subsubsection{배정 --- 어느 방어정이 어느 무리를 맡는가}

적선은 매 결정마다 방위 간격으로 군집화되어 최대 $C$개의 무리가 된다($C=4$). 각 무리 $j$의
위협도는 척수와 모선 근접도의 곱으로 정의한다.
\begin{equation}
\tau_j = n_j \cdot \mathrm{clip}\!\left(1 - \frac{\lVert \mathbf{c}_j - \mathbf{m}\rVert}{R_w},
         \, 0.05,\, 1\right)
\label{eq:threat}
\end{equation}
여기서 $n_j$는 무리의 척수, $\mathbf{c}_j$는 무리 중심, $\mathbf{m}$은 모선 위치,
$R_w$는 교전 해역 반폭이다. 즉 \emph{수가 많고 가까울수록} 위협이 크다.

위협 상위 $P$개 무리를 방어 대상으로 삼고(방어정 수와 같게 하여 1:1), 각 무리의 요격점
$\mathbf{I}_j$를 무리 중심과 모선을 잇는 선분 위에 둔다.
\begin{equation}
\mathbf{I}_j = \mathbf{c}_j + \lambda_{\mathrm{I}}\,(\mathbf{m} - \mathbf{c}_j), \qquad \lambda_{\mathrm{I}} = 0.55
\label{eq:intercept}
\end{equation}
즉 요격점은 무리와 모선 사이를 $0.55 : 0.45$ 로 내분하는 점이며, 모선 쪽으로 약간 치우쳐
방어정이 적보다 느려도 먼저 도달할 수 있게 한다. 방어정 $p$를 무리 $j$에 보내는 비용은
요격점까지의 거리 $d_{pj} = \lVert \mathbf{a}_p - \mathbf{I}_j \rVert$ 로 정의한다.
이 비용으로 탐욕적 1:1 배정을 수행하되, \textbf{직전 배정에 보너스를 준다}. 직전 무리가
아직 방어 대상이면 그 쌍의 비용에서 상수 $\beta$를 뺀다.
\begin{equation}
\tilde{d}_{pj} = d_{pj} - \beta \cdot \mathbb{1}[\,j = j_p^{\text{prev}}\,]
\label{eq:sticky}
\end{equation}

\textbf{남는 방어정의 처리.} 활성 무리 수가 방어정 수보다 적으면 1:1 매칭 뒤 배가
남는다. 이 배들을 놀리지 않고 \emph{가장 과부하인} 무리(부하 $=$ 적 수 $/(1+$배정 수$)$)에
덧붙이되, 같은 자리에 겹쳐 세우지 않는다. $n$번째로 덧붙는 배는 요격점을 모선 쪽으로
$\lambda_{\mathrm{L}} = 0.18$ 만큼 밀어 \emph{한 층 안쪽}에 세우고, 방위로도 $\pm 40^\circ$씩 번갈아
틀어 옆으로 뗀다. 층만 나누고 방위를 두면 두 배가 같은 방위선 위를 앞뒤로 달려 경로가
겹치기 때문이다. 이 규칙이 실제로 담당하는 몫은 작지 않다 --- 활성 무리가 하나뿐인 집중
포메이션에서는 배정의 과반이 이 경로로 결정된다.

이 보너스 항이 이 알고리즘의 핵심이다. 순수하게 거리만 최소화하면 배가 조금 움직일 때마다
최근접 무리가 바뀌어 담당이 요동친다. 방어정이 그물벽을 절반쯤 전개한 상태에서 담당이
바뀌면 그때까지의 이동과 그물이 회수되지 않으므로, \emph{순간 최적을 조금 포기하고 연속성을
사는} 것이 유리하다. \S\ref{sec:res_mechanism}에서 보이듯 이 전환 비용은 언어모델 지휘관의
성능을 가르는 요인이기도 하다.

\subsubsection{경유점 --- 해면 어디에 그물벽을 놓는가}

%-------------------------------------------------- 기호표
%  ★ 이 표에 기호를 추가하려면 본문에 실제로 그 기호가 있는지 먼저 확인할 것.
%    표에만 있고 본문에 없는 기호는 독자를 찾아 헤매게 만든다.
\begin{table*}[t]
\caption{주요 기호. 혼동되기 쉬운 수량은 서로 다른 문자를 쓴다 --- 특히 무리 수 $C$,
         경유점 수 $K$, 학습의 그룹 후보 수 $G$ 는 구분한다.}
\label{tab:nomenclature}
% 다른 표와 같은 \footnotesize 로 둔다. \small(10.95pt)로 두면 자연 폭이 커져
% \resizebox 가 발동하고, 그러면 글자가 오히려 더 작아진다(측정: 10.95 -> 9.53).
\centering\footnotesize
% 판면보다 넓을 때만 줄인다(2단 양단걸침 ≈484pt vs elsarticle preprint 390pt).
\resizebox{\ifdim\width>\linewidth\linewidth\else\width\fi}{!}{%
\begin{tabular}{@{}ll@{\hspace{1.2em}}ll@{}}
\toprule
\textbf{기호} & \textbf{뜻 (값)} & \textbf{기호} & \textbf{뜻 (값)} \\
\midrule
$M$            & 적 USV 수 (10)                       & $\mathbf{o}_p$ & 방어정 $p$ 의 자체 상태 ($\mathbb{R}^8$) \\
$P$            & 방어정 수 (3)                        & $\mathbf{q}_p$ & 방어정 $p$ 의 질의 벡터 ($\mathbb{R}^{32}$) \\
$C$            & 무리 수, 최대 (4)                    & $\mathbf{f}_i$ & 픽셀 $i$ 의 키 벡터 ($\mathbb{R}^{32}$) \\
$K$            & 방어정 1척이 지목하는 경유점 수 (2)  & $s_{p,i}$      & 방어정 $p$ 가 픽셀 $i$ 에 준 점수 \\
$G$            & 학습의 그룹 후보 수 (8)              & $(\boldsymbol{\gamma}_{\mathrm{F}},\boldsymbol{\beta}_{\mathrm{F}})$ & FiLM 배율·이동 계수 \\
\addlinespace[2pt]
$\tau_j$       & 무리 $j$ 의 위협도                   & $R_k$          & 후보 $k$ 의 팀 보상 \\
$d_{pj}$       & 방어정 $p$ $\to$ 무리 $j$ 요격점 거리 & $A_k$         & 후보 $k$ 의 그룹 상대 이득 \\
$\beta$        & 배정 유지 보너스 (식~\ref{eq:sticky}) & $\Phi$        & 위협 근접도 퍼텐셜 \\
$\lambda_{\mathrm{L}}$ & 층 분리 증분 (0.18)                  & $\gamma$       & 할인 계수 (0.99) \\
\addlinespace[2pt]
$\mathbf{I}$   & 요격점                               & $\mathcal{V}^{(k)}_p$ & $k$ 번째 선택 시점의 후보 픽셀 집합 \\
$\mathbf{t}_k$ & $k$ 번째 목표점 (연속 좌표)          & $r_{\mathrm{dup}}$ & 중복 배제 반경 (200\,m) \\
$\hat{\mathbf{r}}$ & 모선 $\to$ 요격점 단위 벡터      & $L_{\mathrm{net}}$ & 그물벽 최대 길이 (450\,m) \\
$\mathbf{a}_p$ & 방어정 $p$ 의 위치                   & $\dz$          & 대응표본 효과크기 \\
\addlinespace[2pt]
$E$, $\ell$    & 관측 반폭 (6300\,m), 픽셀 크기 (252\,m) & $H$           & 후보 평가 창 (180\,step) \\
$\iota(\cdot)$ & 월드 $\to$ 픽셀 인덱스 사상 (식~\ref{eq:world2pix}) & $\eta_t$, $\lambda_{\mathrm{H}}(t)$ & 학습률·엔트로피 계수 스케줄 \\
$\mathbf{u}_k$, $\boldsymbol{\sigma}_o$ & 부화소 오프셋과 그 표준편차 & $w_n$ & 유효 월드 마스크 \\
\bottomrule
\end{tabular}
}
\\[3pt]
\begin{minipage}{\linewidth}\footnotesize
격자는 $50\times50$ 이고 1\,px $=252\unit{m}$ 이다(관측 반폭 6300\,m). 결정 주기는
25\,step $=25$초이며 저수준 제어는 1초마다 돈다. 설정값 전체는
Table~\ref{tab:scenario}에 있다.
\end{minipage}
\end{table*}

\textbf{선행 정의.} 이 절과 다음 절은 같은 이산 행동공간 위에서 동작한다 --- 모선을 중앙에
둔 $50\times50$ 격자(1\,px $=252\unit{m}$)와, 그 격자에서 그물벽을 놓을 수 있는 칸만 남긴
\emph{유효 마스크}이다. 두 대상의 구성 규칙은 \S\ref{sec:obs}과 \S\ref{sec:validmask}에서
정의하며, 여기서는 ``휴리스틱도 정책과 같은 격자·같은 마스크 위에서 점을 고른다''는 사실만
쓴다. 이 공유가 \S\ref{sec:res_maneuver}의 비교를 성립시키는 조건이다.

배정이 확정되면 방어정 $p$는 요격점 $\mathbf{I}$와 그 주변의 유효 영역을 받는다. 휴리스틱은
\textbf{적의 접근 방향에 수직인 벽}을 친다. 모선에서 요격점으로 향하는 단위 벡터를
$\hat{\mathbf{r}}$이라 하면 회랑에 수직인 방향은 $\hat{\mathbf{r}}^{\perp}$이고, $K$개의
목표점을 그 방향으로 등간격 배치한다.
\begin{equation}
\mathbf{t}_k = \mathbf{I} + \left(k - \frac{K-1}{2}\right) L_{\mathrm{net}}\,
               \hat{\mathbf{r}}^{\perp}, \qquad k = 0,\dots,K-1
\label{eq:heur_target}
\end{equation}
$K=2$, $L_{\mathrm{net}} = 450\unit{m}$ 이므로 두 점은 요격점을 가운데 두고 회랑을 가로지르는
길이 450\,m의 선분을 이룬다. 이하에서 $K$ 는 \emph{한 방어정이 찍는 경유점 수}만을 뜻한다 ---
무리 수는 $C$, 학습의 그룹 후보 수는 $G$ 로 구분하여 적는다(Table~\ref{tab:nomenclature}).

$\mathbf{t}_k$는 연속 좌표이므로 그대로 쓸 수 없다. \textbf{유효 픽셀 중 가장 가까운 칸으로
스냅}한다.
\begin{equation}
i_k = \arg\min_{i \in \mathcal{V}^{(k)}_p} \lVert \mathbf{x}_i - \mathbf{t}_k \rVert^2
\end{equation}
여기서 $\mathcal{V}^{(k)}_p$는 방어정 $p$ 의 $k$번째 선택 시점의 후보 집합이고, $\mathbf{x}_i$는 픽셀 $i$의 중심
좌표다. 한 점을 고를 때마다 후보 집합이 갱신된다 --- 반경 $r_{\mathrm{dup}} = 200\unit{m}$
이내를 배제하고(두 점이 겹치면 벽이 서지 않는다), 직전 선택으로부터 $1.5 L_{\mathrm{net}}$
밖도 배제한다(칠 수 없는 길이의 벽을 막는다). 이 배제로 후보가 모두 사라지면 제약을
완화하여 결정을 반환한다.

\textbf{★ 이 스냅과 배제 규칙은 \S\ref{sec:policy}의 학습 정책이 쓰는 것과 같다.}
두 방법이 같은 격자, 같은 유효 마스크, 같은 $r_{\mathrm{dup}}$, 같은 도달성 상한을 통과하므로
\emph{행동공간이 동일하다}. 이것이 \S\ref{sec:res_maneuver}의 비교가 성립하기 위한 전제다.
행동공간이 다르면 성능 차이가 방법의 차이인지 표현력의 차이인지 구분되지 않는다.

\textbf{다만 제약이 모두 소진된 뒤의 처리는 같지 않다.} 휴리스틱은 후보가 비면 그 방어정의
유효 마스크 전체로 한 번에 되돌린다 --- 중복 배제와 도달성 상한을 동시에 푸는 것이다. 정책은
단계별로 푼다. 도달성 상한 때문에 비면 상한만, 다른 방어정의 선점 때문에 비면 선점만 풀고,
그래도 비는 극단에서는 \emph{격자 전체를 유효로 간주}한다. 마지막 단계는 육지나 요격 환형
밖도 후보로 들어올 수 있다는 뜻이므로 휴리스틱에는 대응물이 없다. 이 단계가 실제로 발동한
빈도는 계측하지 않았다. 다만 이 차이는 유효 마스크가 붕괴한 드문 상태에서만 나타나며, 정상
상태의 행동공간이 동일하다는 위 전제를 바꾸지는 않는다.

\subsubsection{이 방법이 표현하지 못하는 것}

휴리스틱은 유효 영역 안에서 \emph{기하적으로 정해진 한 점}을 지목한다. 바꾸어 말하면
유효 영역의 나머지 칸에 대해 ``여기는 저기보다 낫다''는 판단이 없다
(Fig.~\ref{fig:heurvsunet}의 가운데 패널). 다음 요소들이 반영되지 않는다.

\begin{itemize}[leftmargin=*, itemsep=2pt, topsep=3pt]
\item \textbf{다른 무리를 맡은 방어정의 그물}: 앞 문단의 층 분리는 \emph{같은 무리에
      덧붙는 배들 사이}에서만 작동한다. 다른 무리를 맡은 방어정이 이미 친 벽이나 그 벽이
      만들어 낸 차단 상태는 요격점 계산에 들어가지 않는다.
\item \textbf{다음 파의 도착}: 파상 공격에서 앞 단을 막은 벽이 뒤 단까지 막을 수 있는
      자리가 있으나, 요격점은 현재 무리만으로 결정된다.
\item \textbf{육지와 해안 지형}: 유효 마스크가 육지를 배제하기는 하나, 지형이 만드는
      좁은 길목을 \emph{활용}하지는 못한다.
\end{itemize}

\begin{figure*}[t]
\centering
\includegraphics[width=\textwidth]{fig_heur_vs_unet.png}
\caption{기동 계층 두 방식의 비교(학습된 체크포인트의 실제 출력, 같은 전장·같은 시점).
         왼쪽: 두 방식이 공유하는 유효 마스크 --- 2500칸 중 203칸만 선택 가능하다.
         가운데: 휴리스틱은 유효 영역 안에서 기하로 정해진 두 점을 지목하며, 나머지
         칸에 대한 선호를 표현하지 않는다(영역 전체가 균일하게 칠해져 있다).
         오른쪽: 점수맵 정책은 같은 영역 위에 학습된 분포를 얹어 소수의 칸에 확률을
         집중시킨다. \textbf{두 방식은 같은 격자·같은 마스크·같은 거리 제약을
         통과한다} --- 행동공간이 동일해야 성능 차이가 방법의 차이가 된다.}
\label{fig:heurvsunet}
\end{figure*}

이 셋은 모두 ``해면 전체를 놓고 어디가 더 좋은지 비교''해야 얻어지는 판단이다.
\S\ref{sec:policy}의 점수맵은 정확히 그 비교를 학습된 함수로 대체한다.

\subsection{기동 계층 (2): 점수맵 위 픽셀 지목 정책}
\label{sec:policy}

\subsubsection{다채널 격자 관측}
\label{sec:obs}

기동 계층의 입력은 좌표 목록이 아니라, 모선을 중앙에 두고 전장을 $50\times50$ 격자로 나눈
다채널 지도이다. 채널 하나하나가 실제로 어떤 그림인지는 Fig.~\ref{fig:observation}에 실었다. 픽셀 한 변은 252\,m 이고 관측 반폭은
6300\,m 이다. 총 15개 채널은 세 묶음으로 구성된다(Fig.~\ref{fig:tensorstack}).

\begin{itemize}[leftmargin=*, itemsep=2pt, topsep=3pt]
\item \textbf{전역 9채널} (모든 방어정이 공유): 적 밀도, 적 진행 $x$, 적 진행 $y$, 적 위협도,
      아군 위치, 설치된 그물, 원점, 요격 환형, 육지.
\item \textbf{방어정별 3채널}: 자기 위치, 배정 요격점, 유효 영역.
\item \textbf{좌표 3채널}: CoordConv~\citep{liu2018coordconv} 의 $x$, $y$, $r$.
\end{itemize}

\textbf{래스터화와 정규화.} 격자 한 변을 $n_{\mathrm{g}}=50$, 관측 반폭을 $E = 6300\unit{m}$ 라 하면
픽셀 크기는 $\ell = 2E/n_{\mathrm{g}} = 252\unit{m}$ 이고, 월드 좌표 $\mathbf{x}$ 는 모선 위치
$\mathbf{m}$ 을 중심으로 다음 픽셀 인덱스로 사상된다.
\begin{equation}
\iota(\mathbf{x}) = \mathrm{clip}\!\left(\left\lfloor \frac{\mathbf{x} - \mathbf{m} + E\mathbf{1}}{\ell} \right\rfloor,\, 0,\, n_{\mathrm{g}}-1\right)
\label{eq:world2pix}
\end{equation}
적 관련 채널은 개체를 픽셀에 ``칠하는'' 방식으로 만들며, 값 범위가 $[0,1]$ 또는 $[-1,1]$
에 놓이도록 정규화한다. 생존 적선 집합을 $\mathcal{A}$, 적선 $m$ 의 위치·헤딩을
$\mathbf{e}_m$, $\psi_m$ 이라 하면 픽셀 $i$ 의 적 밀도·진행 방향·위협도 채널은
\begin{equation}
c^{\mathrm{den}}_i = \frac{N_i}{M}, \quad
c^{\mathrm{vel}}_i = \frac{1}{\max(N_i, 1)} \sum_{m \in \mathcal{A}_i} (\sin\psi_m, \cos\psi_m)
\label{eq:ch_enemy}
\end{equation}
\begin{equation}
c^{\mathrm{thr}}_i = \max_{m \in \mathcal{A}_i}
   \left(1 - \mathrm{clip}\!\left(\frac{\lVert \mathbf{e}_m - \mathbf{m} \rVert}{v_e T_{\max}},\, 0,\, 1\right)\right)
\label{eq:ch_threat}
\end{equation}
이다. 여기서 $\mathcal{A}_i = \{m \in \mathcal{A} : \iota(\mathbf{e}_m) = i\}$ 는 픽셀 $i$ 에
놓인 생존 적선의 집합이고 $N_i = |\mathcal{A}_i|$ 이다. 밀도는 적 총수 $M$ 으로, 진행 방향은 픽셀 안의 적 수로 나누어 평균 단위벡터가 되게
하고, 위협도는 모선까지 거리를 적이 에피소드 상한 안에 달릴 수 있는 최대 거리
$v_e T_{\max}$ ($9\unit{m/s} \times 2000\unit{s}$)로 정규화하여 가까울수록 1에 가깝게 한다.
아군 위치 채널은 같은 방식으로 $P$ 로 나누고, 설치된 그물 채널은 고해상도 물리 격자의
최대값 풀링, 육지 채널은 면적비 다운샘플(Fig.~\ref{fig:landpipe})이다. 좌표 채널은
픽셀 중심 $\mathbf{x}_i = (x_i, y_i)$ 에 대해
\begin{equation}
(c^{x}_i,\, c^{y}_i) = \frac{\mathbf{x}_i - \mathbf{m}}{E}, \qquad
c^{r}_i = \mathrm{clip}\!\left(\frac{\lVert \mathbf{x}_i - \mathbf{m} \rVert}{E},\, 0,\, 1\right)
\label{eq:ch_coord}
\end{equation}
로 두어 $[-1,1]$ 과 $[0,1]$ 에 놓는다. 격자에 담지 않는 방어정 $p$ 의 자체 상태 벡터는
같은 반폭 $E$ 로 정규화한 위치와 요격점, 헤딩 단위벡터, 그물 잔량 비율, 전개 여부이다.
\begin{equation}
\begin{aligned}
\mathbf{o}_p = \Big[&\tfrac{1}{E}(\mathbf{a}_p - \mathbf{m});\; (\sin\psi_p, \cos\psi_p);\;
  \tfrac{n_p}{n_{\max}};\; \delta_p;\\
  &\mathbb{1}[\text{배정}]\,\tfrac{1}{E}(\mathbf{I}_p - \mathbf{m})\Big] \in \mathbb{R}^{8}
\end{aligned}
\label{eq:own}
\end{equation}
여기서 $n_p$ 는 잔여 그물 수, $n_{\max}$ 는 척당 보유량, $\delta_p \in \{0,1\}$ 은 전개 중
여부이며, 배정이 없으면 요격점 성분은 0이다. 이렇게 모든 입력이 $[-1,1]$ 안에 들어오므로
별도의 통계 기반 정규화 층이 필요 없고, 학습 분포와 배치 분포의 정규화 불일치가 생기지
않는다.

\begin{figure*}[t]
\centering
\includegraphics[width=0.96\textwidth]{fig_tensor_stack.png}
\caption{15채널 관측의 구성. 전역 9채널(01--09), 방어정별 3채널(10--12), 좌표 3채널
         (13--15)이 같은 $50\times50$ 격자에 겹쳐 하나의 다채널 영상을 이룬다. 오른쪽의
         채널 이름은 구현의 이름과 같다. 왼쪽 아래는 격자에 담기지 않고 벡터로 들어가는
         방어정 자체 상태 $\mathbf{o}_p \in \mathbb{R}^8$ --- 위치 2, 헤딩 단위벡터 2,
         잔여 그물 1, 전개 중 여부 1, 배정 요격점 2이며, 요격점 두 성분은 배정이 없으면
         0이다. \emph{그림의 좌측 시선원뿔(``egocentric FOV'')과 함정 형상은 삽화의
         표현이다} --- 본 연구의 래스터는 방어정 시야가 아니라 \textbf{모선을 중앙에 둔
         전역 격자}이며, 방어정마다 달라지는 것은 3채널(10--12)과 $\mathbf{o}_p$ 뿐이다.}
\label{fig:tensorstack}
\end{figure*}

Fig.~\ref{fig:observation}은 이 15채널 관측을 한 장면에 대해 채널별로 펼친 것이다.
채널의 값 범위와 비영 픽셀 수를 함께 적어, 어떤 채널이 실제로 정보를 담고 어떤 채널이
대부분 0인지가 드러나게 하였다 --- 예컨대 좌표 채널은 전 픽셀이 비영이고, 방어정별 채널은
수십 픽셀 규모이며, 설치된 그물 채널은 전개 이전이라 전부 0이다.

\begin{figure*}[t]
\centering
\includegraphics[width=\linewidth]{fig_channels.png}
\caption{래스터화된 다채널 관측의 채널별 분해. 전역 채널(적 밀도, 적 진행 $x$·$y$,
         적 위협도, 아군 위치, 설치된 그물, 원점, 요격 환형), 방어정별 채널(자기 위치,
         배정 요격점, 유효 영역), 좌표 채널($x$, $y$, $r$)이 같은 격자에 정렬된다.
         각 패널에 min/max와 비영 픽셀 수를 함께 표시하였다.
         \emph{이 그림은 시각적 판독을 위해 $100\times100$ (126\,m/px)로 렌더하였다} ---
         본문의 학습·평가 설정은 $50\times50$ (252\,m/px)이며 채널 구성과 의미는 동일하다.
         육지 채널은 이 장면에 육지가 없어 표시되지 않았다.}
\label{fig:observation}
\end{figure*}

목록 대신 격자 지도를 쓰면 두 가지 이득이 있다. 교전 중 적선 수가 늘거나 줄어도 입력 형태가
불변이고, 적선을 어떤 순서로 넣었는지에 따라 결과가 바뀌지 않는다(순열 불변). 다만 합성곱은
평행이동 등변이라 절대 위치가 흐려지므로, 좌표 채널을 함께 넣어 위치를 보존한다.

육지 채널은 실제 해역 정보를 반영한다(Fig.~\ref{fig:landpipe}). OpenStreetMap 에서 해안선·섬 폴리곤을 추출하여
격자로 변환한 뒤, 열림 연산으로 부두·방파제 같은 가는 돌출부를 제거하고 구멍 메우기로
내부를 정리한다. 충돌 판정은 고해상도 물리 격자에서 이진으로 수행하고, 관측 채널은 이를
면적비로 $50\times50$ 에 다운샘플하여 해안선이 계단이 아닌 연속값으로 들어가게 한다.
해역별로 한 번만 계산하여 캐시하므로 학습 중 외부 API 호출은 발생하지 않는다.

\begin{figure*}[t]
\centering
\includegraphics[width=\textwidth]{fig_landpipe.png}
\caption{실해역 지형의 반영 과정(통영 매물도 인근, 12.6\,km 박스). ① OpenStreetMap
         벡터 폴리곤 $\to$ ② 고해상도 래스터($805^2$, 16\,m/px) $\to$ ③ 충돌 판정용
         물리 격자($200^2$, 63\,m/cell) $\to$ ④ 정책 격자($50^2$, 252\,m/px)로 면적비
         다운샘플. \textbf{판정과 관측을 다른 해상도로 두는 이유}는 충돌은 이진 판정을
         요구하고, 관측은 해안선이 계단으로 들어가면 안 되기 때문이다. ⑤는 적의 우회를
         유도하는 척력장으로, 관측이 아니라 적 기동 모형의 일부다.}
\label{fig:landpipe}
\end{figure*}

%-------------------------------------------------- 그림 4

\subsubsection{FiLM 조건화 U-Net과 점수맵}


관측은 U-Net 구조~\citep{ronneberger2015unet}를 지나 입력과 같은 해상도의 특징 지도로
변환된다(Fig.~\ref{fig:policy}). 인코더 2단($50\to25\to13$)으로 수용 영역을 넓히고, 디코더
2단으로 원해상도를 복원하며, skip 연결이 인코더의 위치 정보를 디코더로 전달한다. 정밀 좌표와
전역 상황 판단을 동시에 요구하는 본 문제에 U자형 구조가 맞물리는 지점이 여기다.

병목에서는 전역 평균 풀링으로 압축한 전황 벡터로부터 FiLM~\citep{perez2018film} 의
채널별 배율·이동 계수 $(\boldsymbol{\gamma}_{\mathrm{F}}, \boldsymbol{\beta}_{\mathrm{F}})$ 를
만들어 병목 특징을 변조한 뒤 디코더에 넘긴다(아래첨자 F 는 FiLM 을 뜻하며, 식~\eqref{eq:sticky}의 $\beta$ 나
할인 계수 $\gamma$ 와는 무관하다). 병목 특징 지도를 $\mathbf{H}_b \in \mathbb{R}^{2w \times h \times h}$
($w = 32$, $h = 13$)라 하면
\begin{equation}
\begin{aligned}
(\boldsymbol{\gamma}_{\mathrm{F}}, \boldsymbol{\beta}_{\mathrm{F}})
  &= \mathbf{W}_{\mathrm{F}}\, \mathrm{GAP}(\mathbf{H}_b) + \mathbf{b}_{\mathrm{F}},\\
\tilde{\mathbf{H}}_b &= (1 + \boldsymbol{\gamma}_{\mathrm{F}}) \odot \mathbf{H}_b + \boldsymbol{\beta}_{\mathrm{F}}
\end{aligned}
\label{eq:film}
\end{equation}
이며, $\odot$ 은 채널별 곱이다. 배율을 $1 + \boldsymbol{\gamma}_{\mathrm{F}}$ 로 두어 초기에는
항등 변조에서 출발한다. 국소 합성곱만으로는 ``지금 전황이 급박한가''와 같은 전역 정보가
모든 위치에 고르게 전달되지 않기 때문에 이 경로를 둔다.

각 방어정 $p$ 는 자신의 상태 벡터 $\mathbf{o}_p \in \mathbb{R}^8$(식~\eqref{eq:own})로부터
질의 벡터 $\mathbf{q}_p = \mathrm{MLP}_{8 \to 128 \to 32}(\mathbf{o}_p) \in \mathbb{R}^{32}$ 를
만들고, 픽셀 $i$ 의 특징 $\mathbf{f}_i \in \mathbb{R}^{32}$ 와의 내적으로 점수를 얻는다.
\begin{equation}
s_{p,i} = \frac{\langle \mathbf{W}_q \mathbf{q}_p + \mathbf{b}_q,\, \mathbf{f}_i \rangle}{\sqrt{d}}, \qquad d = 32
\label{eq:score}
\end{equation}
여기서 $\mathbf{W}_q \in \mathbb{R}^{d \times d}$, $\mathbf{b}_q \in \mathbb{R}^{d}$ 는 질의를 특징 공간에 맞추는 아핀 사영이다.
이는 어텐션~\citep{vaswani2017attention}의 점수 계산과 같은 형식이며, 포인터
방식~\citep{vinyals2015pointer}으로 격자 위의 원소를 지목하는 구조에 해당한다. 특징 지도는
방어정마다 다르므로(15채널 중 3채널이 방어정별이다), 관측 텐서는 방어정을 배치 축으로
쌓아 \textbf{결정당 한 번의 배치 호출}로 U-Net을 통과한다. 합성곱 연산량은 방어정 수에
비례하지만, \textbf{가중치는 모든 방어정이 공유하므로 모델 크기는 방어정 수에 불변}이다.
방어정별로 달라지는 것은 질의 $\mathbf{q}_p$ 와 방어정별 3채널뿐이다.

합성곱 블록은 $3\times3$ 합성곱 --- 그룹 정규화~\citep{wu2018groupnorm} --- SiLU 로 구성한다.
\begin{equation}
\begin{aligned}
\mathrm{Block}(\mathbf{H}) &= \mathrm{SiLU}\big(\mathrm{GN}_{8}(\mathrm{Conv}_{3\times3}(\mathbf{H}))\big),\\
\mathrm{SiLU}(z) &= z / (1 + e^{-z})
\end{aligned}
\label{eq:block}
\end{equation}
그룹 정규화($\mathrm{GN}_8$, 8개 그룹)를 고른 이유는 성능이 아니라 배치 독립성이다. 표본
하나 안에서 채널을 묶어 정규화하므로 통계가 배치 구성에 의존하지 않으며, 이것이
\S\ref{sec:deploy}에서 다루는 배치 크기 1의 실시간 추론이 학습과 같은 연산을 수행한다는
전제가 된다.

%-------------------------------------------------- 그림 5 (양단)
\begin{figure*}[t]
\centering
\includegraphics[width=0.94\textwidth]{fig_unet_arch.png}
\caption{점수맵 정책의 구조. 각 방어정의 8차원 상태가 공유 MLP($8\to128\to32$)를
         지나 질의 $\mathbf{q}$가 되고, 15채널 관측은 Stem($32\times50\times50$) --
         인코더($32\times25\times25 \to 64\times13\times13$) -- FiLM 병목 --
         디코더($32\times25\times25 \to 32\times50\times50$)를 지난다. skip 연결이
         인코더의 위치 정보를 디코더로 전달한다. 질의와 픽셀 특징의 내적이
         $50\times50$ 점수맵을 이루고, 회귀 머리가 선택된 칸 안의 부화소 오프셋 2개를
         낸다. 방어정을 배치 축으로 쌓아 \textbf{결정당 한 번의 배치 호출}로 통과하며
         (그림의 입력 부피는 방어정마다 하나씩이다), 가중치는 전 방어정이 공유하므로
         모델 크기는 방어정 수에 불변이다.}
\label{fig:policy}
\end{figure*}

Fig.~\ref{fig:policy}가 구조의 개념 일러스트라면, Fig.~\ref{fig:pipeline_detail}은
같은 망의 형상과 조건화 지점을 정확히 적은 구현 참조도이다 --- 실제 관측 채널 화면에서
출발해 각 층의 텐서 크기, FiLM 계수가 들어가는 위치, 자기회귀 질의 갱신까지 한 장에 담았다.

\begin{figure*}[t]
\centering
\includegraphics[width=\textwidth]{fig_unet_plotnn.png}
\caption{기동 계층의 텐서 흐름(구현 참조도, 244{,}389 파라미터, $d=32$, $K=2$).
         \textbf{왼쪽}: 실제 관측 15채널 --- 전역 9채널(위)과 방어정별 3채널 $+$ CoordConv
         3채널(아래)의 학습된 체크포인트 출력 화면. \textbf{가운데}: stem $\to$ enc1 $\to$
         enc2(stride 2) $\to$ bottleneck $\to$ dec2 $\to$ dec1의 U-Net. 병목에서 전역 평균
         풀링 $\to$ 선형층으로 FiLM 계수 $(\gamma,\beta)$를 만들어 병목 특징을 변조하고,
         방어정 자체 상태 $\mathbf{o}_p \in \mathbb{R}^8$ 은 MLP($8\to128\to32$)를 거쳐
         질의 $q$ 가 된다. 스킵 연결은
         위쪽 파란 화살표. \textbf{오른쪽}: $1\times1$ 합성곱으로 픽셀 특징 $f$를 내고,
         질의 $q$와의 내적 $\langle q, f\rangle/\sqrt{d}$로 점수맵을 만든 뒤 유효 마스크를
         거쳐 argmax로 $K=2$ 픽셀을 고른다. 아래 점선은 자기회귀
         $q \leftarrow q + f_{k-1}$ --- 첫 점의 특징을 질의에 더해 두 번째를 고른다.
         본문 Fig.~\ref{fig:policy}는 같은 구조의 개념 일러스트이고, 이 그림은 형상과
         조건화 지점을 정확히 적은 판이다.}
\label{fig:pipeline_detail}
\end{figure*}

\subsubsection{유효 마스크와 픽셀 지목}
\label{sec:validmask}

점수맵 전체가 후보가 되는 것은 아니다. 다음 조건을 순차적으로 적용하여 유효 마스크를
구성하고, 그 밖은 $-\infty$ 로 차단한다.
\begin{enumerate}[leftmargin=*, itemsep=1pt, topsep=3pt]
\item 요격 가능 환형(모선에서 일정 반경 대역) 안쪽
\item 육지가 아닌 곳
\item 배정된 요격점으로부터 반경 게이트 이내
\item 배정 무리의 방위 회랑 이내
\end{enumerate}
픽셀 $i$ 의 중심을 $\mathbf{x}_i$, 육지 픽셀 집합을 $\mathcal{L}_{\mathrm{land}}$, 방어정 $p$ 의 배정
요격점을 $\mathbf{I}_p$, 모선에서 본 요격점 방위를 $\vartheta_p$, 픽셀 방위를 $\vartheta_i$ 라
하면 방어정 $p$ 의 유효 픽셀 집합은 네 조건의 교집합이다.
\begin{equation}
\begin{aligned}
\mathcal{V}_p = \big\{\, i :\; & r_{\min} \le \lVert \mathbf{x}_i - \mathbf{m} \rVert \le r_{\max},\;
  i \notin \mathcal{L}_{\mathrm{land}},\\
  & \lVert \mathbf{x}_i - \mathbf{I}_p \rVert \le r_{\mathrm{g}},\;
  |\vartheta_i - \vartheta_p|_{\circ} \le \vartheta_{\mathrm{g}} \,\big\}
\end{aligned}
\label{eq:validset}
\end{equation}
여기서 $|\cdot|_{\circ}$ 는 $[0^\circ, 180^\circ]$ 로 접은 각도 차이이며 $r_{\min} = 400$,
$r_{\max} = 4500$, $r_{\mathrm{g}} = 3000\unit{m}$, $\vartheta_{\mathrm{g}} = 60^\circ$ 이다.
유효 픽셀이 하한 150개에 못 미치면 반경 게이트부터 푼다 --- 방위 회랑은 ``무엇을 막는
그물인가''를 정하므로 마지막까지 유지한다. 마스크는 점수에 곱해지는 것이 아니라 로짓
단계에서 적용된다.
\begin{equation}
\tilde{s}_{p,i} = \begin{cases} s_{p,i}, & i \in \mathcal{V}_p \\ -\infty, & i \notin \mathcal{V}_p, \end{cases}
\quad
\pi_\theta(i \mid \mathbf{s}, p) = \frac{e^{\tilde{s}_{p,i}}}{\sum_{j \in \mathcal{V}_p} e^{\tilde{s}_{p,j}}}
\label{eq:masked_softmax}
\end{equation}
유효 집합 밖 픽셀의 확률은 정확히 0 이다. 이 마스킹은 안전·임무 제약을 \emph{행동 분포에 직접} 부과한다는 점에서 중요하다. 연속 좌표를
회귀시키는 방식에서는 이런 제약을 손실 항으로 넣어 ``되도록 지키게'' 유도할 수밖에 없지만,
지목 방식에서는 제약 위반이 애초에 표현 불가능하다.

이 제약이 실제로 얼마나 강한지는 측정으로 확인된다(Fig.~\ref{fig:validmask}). 학습된
체크포인트로 6{,}652회의 결정을 수집한 결과, $50\times50=2500$개 칸 가운데 평균
\textbf{10.5\,\%}만이 선택 가능하였다. 대표 사례의 단계별 감소는 전체 2500칸 $\to$ 요격
환형 1000 $\to$ 육지 제외 979 $\to$ 반경 게이트 387 $\to$ 방위 회랑 299로, 격자의
\textbf{88\,\%가 애초에 갈 수 없는 자리}이다. 연속 좌표 회귀에는 이 구조를 표현할 수단이
없다.

\textbf{왜 연속 좌표가 아니라 이산 격자인가.} 마스킹 말고도 세 가지 이득이 있다.
첫째 \emph{다봉 분포}이다. 그물벽을 놓을 만한 자리는 흔히 한 곳이 아니라 서로 떨어진
여러 곳이다. 연속 좌표에 흔히 쓰는 단봉 가우시안 정책은 이런 후보들을 평균 내어 그 사이의
--- 종종 유효하지도 않은 --- 지점으로 수렴하지만, 격자 위 범주형 분포는 떨어진 봉우리를
그대로 보존한다. 둘째 \emph{탐색의 유계성}이다. 범주형 표본은 정의상 유효 영역을 벗어나지
않으므로 좌표를 자르거나 squash 하는 장치가 필요 없고, 엔트로피가 유한하여
\S\ref{sec:training}의 그룹 후보들이 서로 다른 칸을 실제로 밟는다 --- 가우시안 정책은
표준편차가 줄면 후보가 한 점으로 뭉쳐 상대 이득(식~\eqref{eq:adv})이 사라진다. 셋째
\emph{구조의 일치}이다. 합성곱의 평행이동 등변성은 픽셀 단위 출력에서 그대로 살아나므로,
관측 격자와 행동 격자를 같은 $50\times50$ 으로 두면 ``이 칸이 좋다''는 판단이 국소
기하로부터 곧바로 나온다. 대가는 해상도이다. 한 칸이 252\,m 이므로 순수 격자 행동은 최대
$\pm126\,\mathrm{m}$ 의 양자화 오차를 갖는데, 이를 위해 선택된 칸 안에서 $\pm$반픽셀의
연속 오프셋을 내는 회귀 머리를 함께 둔다(Fig.~\ref{fig:policy}). 즉 \emph{어느 칸인가}는
이산으로, \emph{칸 안 어디인가}는 연속으로 결정하여 두 표현의 장점을 나누어 갖는다.

%-------------------------------------------------- 유효 마스크
\begin{figure*}[t]
\centering
\includegraphics[width=\linewidth]{fig_valid_mask.png}
\caption{왜 좌표 회귀가 아니라 점수맵 위 픽셀 지목인가(학습된 체크포인트의 실제 출력).
         \textbf{왼쪽}: 연속 좌표를 회귀하는 출력에는 ``놓을 수 없는 곳''을 표현할 자리가 없다
         --- 이 장면에서 붉은 영역이 격자의 88.0\,\%다. \textbf{오른쪽}: 같은 장면을 후보
         영역으로 확대한 점수 분포와 선택된 두 픽셀, 그리고 그 둘을 잇는 그물벽이다.
         마스킹된 softmax 아래에서 부적합한 칸은 확률이 0이라 애초에 선택되지 않는다.
         \textbf{아래 막대}는 제약이 순차 적용되며 후보가 줄어드는 과정이다
         ($2500 \to 1000 \to 979 \to 387 \to 299$, 12.0\,\%). 전 결정 평균 후보는
         10.5\,\%였다.}
\label{fig:validmask}
\end{figure*}

각 방어정은 유효 영역 안에서 점수 상위 픽셀 $K=2$ 개를 \emph{순차적으로} 고른다. $k$ 번째
선택의 질의와 후보 집합을 $\mathbf{q}^{(k)}_p$, $\mathcal{V}^{(k)}_p$ 라 하면
($\mathbf{q}^{(0)}_p = \mathbf{q}_p$, $\mathcal{V}^{(0)}_p = \mathcal{V}_p$)
\begin{equation}
i_k \sim \pi^{(k)}_p \equiv \pi_\theta(\cdot \mid \mathbf{q}^{(k)}_p, \mathcal{V}^{(k)}_p), \qquad
\mathbf{q}^{(k+1)}_p = \mathbf{q}^{(k)}_p + \mathbf{f}_{i_k}
\label{eq:autoreg}
\end{equation}
\begin{equation}
\mathcal{V}^{(k+1)}_p = \big\{ i \in \mathcal{V}^{(k)}_p :
  r_{\mathrm{dup}} \le \lVert \mathbf{x}_i - \mathbf{x}_{i_k} \rVert \le 1.5\, L_{\mathrm{net}} \big\}
\label{eq:exclude}
\end{equation}
이다. 먼저 고른 픽셀의 특징을 질의에 더해 두 번째를 고르므로 두 점이 서로를 알고 뽑힌다.
독립적으로 뽑으면 두 점을 잇는 벽의 방향이 우연에 맡겨진다. 식~\eqref{eq:exclude}는
\emph{두 점 사이 간격}의 상·하한이다. $r_{\mathrm{dup}} = 200\unit{m}$ 안쪽은 두 점이 너무
가까워 벽이 서지 않고, 상한은 그물벽 최대 길이의 1.5배인 675\,m 로 두어 실제로 칠 수 없는
길이의 벽이 선택되지 않게 한다. 상한을 두지 않았을 때 휴리스틱 목표점을 유효 픽셀로 스냅한 간격의
중앙값이 504\,m 로 그물 길이 450\,m 를 초과했고(초과율 66.7\,\%), 정책의 표본은 88.6\,\% 가
초과했다. 이 제약으로 후보가 모두 사라지면 제약을
완화하여 결정을 반환한다(최적성보다 결정 가능성을 우선한다).

선택된 칸 안의 위치는 연속 오프셋으로 정한다. 질의와 선택 픽셀의 특징으로부터 평균을 내고,
학습되는 로그표준편차 $\boldsymbol{\sigma}_o$ 를 가진 가우시안에서 표집한다.
\begin{equation}
\begin{aligned}
\boldsymbol{\mu}_k &= \mathrm{MLP}_{2d \to d \to 2}\big([\mathbf{q}^{(0)}_p;\, \mathbf{f}_{i_k}]\big), \quad
\mathbf{u}_k \sim \mathcal{N}(\boldsymbol{\mu}_k, \boldsymbol{\sigma}_o^2),\\
\mathbf{t}_k &= \mathbf{x}_{i_k} + \tfrac{\ell}{2}\,\mathrm{clip}(\mathbf{u}_k, -1, 1)
\end{aligned}
\label{eq:offset}
\end{equation}
즉 \emph{어느 칸인가}는 범주형으로, \emph{칸 안 어디인가}는 가우시안으로 결정된다. 한
방어정의 행동 $a_p = \{(i_k, \mathbf{u}_k)\}_{k=0}^{K-1}$ 의 로그확률과 엔트로피는 두 분포의
합이다.
\begin{equation}
\log \pi_\theta(a_p) = \sum_{k=0}^{K-1} \Big[ \log \pi^{(k)}_p(i_k)
  + \log \mathcal{N}(\mathbf{u}_k; \boldsymbol{\mu}_k, \boldsymbol{\sigma}_o^2) \Big]
\label{eq:logp}
\end{equation}
\begin{equation}
\mathcal{H}(a_p) = \sum_{k=0}^{K-1} \Big[ \mathcal{H}\big(\pi^{(k)}_p\big)
  + \mathcal{H}\big(\mathcal{N}(\boldsymbol{\mu}_k, \boldsymbol{\sigma}_o^2)\big) \Big]
\label{eq:entropy}
\end{equation}
식~\eqref{eq:logp}는 \S\ref{sec:training}의 정책경사에, 식~\eqref{eq:entropy}는 엔트로피
정칙화에 그대로 쓰인다. 특징 조회에서 두 번째 점까지의 흐름을 Fig.~\ref{fig:picksteps}에
요약하였다.

%-------------------------------------------------- 픽셀 지목 4단계
\begin{figure*}[t]
\centering
\includegraphics[width=0.96\textwidth]{fig_pick_steps.png}
\caption{점수맵 위 픽셀 지목의 네 단계. ① 복호기가 낸 픽셀별 특징 $F(x,y)$.
         ② 유효 마스크 $M$을 곱해 놓을 수 없는 칸을 차단한다(\S\ref{sec:validmask}).
         ③ 질의 $\mathbf{v}$와의 내적으로 점수맵 $S = F\cdot\mathbf{v}$를 얻고 최대 칸을
         첫 점으로 고른다. ④ 첫 점의 특징을 질의에 더해($\mathbf{q} \leftarrow
         \mathbf{q} + \mathbf{f}_0$) 같은 절차로 두 번째 점을 고르되, 첫 점과의 간격
         상·하한(파선 원)을 마스크에 추가한다. 두 점을 잇는 선분이 그물벽이다.}
\label{fig:picksteps}
\end{figure*}

선택된 두 좌표를 잇는 선분이 그물벽 경로가 된다. 방어정은 가까운 점까지 무전개로 이동한 뒤
다음 구간에서 그물을 전개한다. 운용 시에는 최대 점수를 취하고(greedy), 학습 시에는 확률
분포에서 표본을 뽑는다.

\textbf{정지 규칙.} 그물이 없거나 배정되지 않은 방어정은 정지시킨다. 이런 방어정에는 학습
중 유의미한 보상 신호가 없어 점수맵이 거의 평평해지는데(최고 확률 0.02 수준, 배정된 방어정은
$0.5\sim0.7$), 이 무의미한 출력이 항해 명령으로 새어 나가는 것을 막기 위함이다.
\subsection{학습: 가치망 없는 다중 에이전트 그룹 상대 최적화}
\label{sec:training}

\subsubsection{그룹 상대 이득}

학습은 가치함수를 따로 두지 않는 그룹 상대 정책 최적화~\citep{shao2024deepseekmath}를 다중
방어정 상황에 맞게 변형하여 수행한다. 결정 시점의 상태를 $\mathbf{s}$, 방어정 $P$척의 결합
행동을 $\mathbf{a} = (a_1, \dots, a_P)$ 라 하자. 매 결정 시점에서 정책으로부터 $G$개의 결합
행동 후보를 표집하되, 후보 가운데 하나($k=0$)는 휴리스틱(\S\ref{sec:heuristic})이 내는
행동으로 고정하여 기준선으로 삼는다.
\begin{equation}
\mathbf{a}^{(0)} = \mathbf{a}_{\mathrm{heur}}(\mathbf{s}), \qquad
\mathbf{a}^{(k)} \sim \pi_\theta(\cdot \mid \mathbf{s}), \quad k = 1, \dots, G-1
\label{eq:candidates}
\end{equation}
각 후보를 시뮬레이터에서 공통 구간 $H = 180\,\mathrm{step}$ 만큼 실제로 진행시켜 창 보상
$R_k = R(\mathbf{s}, \mathbf{a}^{(k)}; H)$ 를 얻은 뒤(\S\ref{sec:reward}), 후보들 사이의
상대 비교로 이득을 산정한다.
\begin{equation}
A_k = \frac{R_k - R_0}{\mathrm{std}(\{R_j\}_{j=0}^{G-1}) + \epsilon}, \qquad \epsilon = 10^{-6}
\label{eq:adv}
\end{equation}
$R_0 = R_{\mathrm{heur}}$ 이므로 휴리스틱과 같은 성과를 내는 후보의 이득은 0 이고, 휴리스틱을
넘어서는 후보만 양의 이득을 받는다. 그룹 표준편차로 나누므로 이득은 크기가 없는 양이며,
그룹 안의 순위 정보만 남는다. 어떤 선택이 더 나았는지를 실제 결과로 직접 비교하므로 별도의
가치망이 필요하지 않다. 이는 가치 추정 편향과 비평가 학습의 불안정성을 원천적으로 제거한다.

한 갱신은 $N$개의 병렬 전장(월드)에서 동시에 수행된다. 월드 $n$ 의 후보 보상이 모두 같으면
이득이 정의되지 않으므로, 유효 마스크 $w_n = \mathbb{1}[\mathrm{std}(\{R_{n,k}\}_k) > \epsilon_{\mathrm{var}}]$
($\epsilon_{\mathrm{var}} = 10^{-6}$)로 그런 월드를 갱신에서 제외한다. 정책경사 손실은 유효 월드의 모든 후보·모든 방어정에 대해 이득으로 가중한 로그확률
(식~\eqref{eq:logp})의 음수이다.
\begin{equation}
\mathcal{L}_{\mathrm{PG}}(\theta) = -\frac{\sum_{n} w_n \sum_{k=0}^{G-1} A_{n,k} \sum_{p=1}^{P} \log \pi_\theta\big(a^{(k)}_{n,p} \mid \mathbf{s}_n\big)}{G \sum_n w_n}
\label{eq:lpg}
\end{equation}
여기에 엔트로피 정칙화(식~\eqref{eq:entropy})를 더한 총손실을 Adam 으로 최소화한다.
\begin{equation}
\mathcal{L}(\theta) = \mathcal{L}_{\mathrm{PG}}(\theta)
  - \lambda_{\mathrm{H}}(t)\, \frac{1}{NP} \sum_{k=0}^{G-1} \sum_{n=1}^{N} \sum_{p=1}^{P} \mathcal{H}\big(a^{(k)}_{n,p}\big)
\label{eq:loss}
\end{equation}
엔트로피는 $G$개 후보 모두에 대해 합하므로 유효 계수는 $G \lambda_{\mathrm{H}}(t)$ 이며,
$\lambda_{\mathrm{H}}(t)$ 는 \S\ref{sec:schedule}의 감쇠 계수이다. 기울기는 노름
$g_{\max} = 1.0$ 에서 자른다.
\begin{equation}
\mathbf{g} = \nabla_\theta \mathcal{L}(\theta), \qquad
\theta \leftarrow \theta - \eta_t \cdot \mathrm{Adam}\!\left(\mathbf{g} \cdot \min\!\left(1, \frac{g_{\max}}{\lVert \mathbf{g} \rVert}\right)\right)
\label{eq:update}
\end{equation} 시뮬레이터를 직접 전진시켜 후보를
평가하는 착상은 탐색 기반 의사결정~\citep{silver2016alphago}과 같은 계열이되, 본 연구는
탐색 트리를 전개하지 않고 한 결정 시점의 후보 집합만 비교한다. 이 절차와 다음 소절의
크레딧 배분을 Fig.~\ref{fig:grpocredit}에 한 장으로 정리하였다.

%-------------------------------------------------- GRPO + 반사실 크레딧
\begin{figure*}[t]
\centering
\includegraphics[width=0.94\textwidth]{fig_grpo_credit.png}
\caption{가치망 없는 다중 방어정 그룹 상대 최적화의 한 갱신. \textbf{위}: 한 결정 상태에서
         $G$개의 결합 행동 후보를 표집하고, 각 후보를 시뮬레이터에서 공통 구간만큼 실제로
         진행시켜 팀 보상을 얻는다. \textbf{가운데}: 후보 평가에 쓰이는 롤아웃 창.
         \textbf{아래 왼쪽·가운데}: 팀 보상을 후보 간 상대 비교로 이득으로 바꾼다
         (식~\eqref{eq:adv}). \textbf{아래 오른쪽}: 반사실적 크레딧 --- 한 척의 행동만 바꾸고
         나머지는 고정한 채 시뮬레이터를 다시 진행시켜, 그 차이를 그 배의 몫으로 돌린다
         (\S\ref{sec:training}; 구현된 변형이며 채택 런은 팀 이득 손실을 썼다). 그림은 기전 전달을 위해 후보 수와 창을 $G=4$, $H=4$로
         줄여 그렸으며 실제 학습은 $G=8$, 평가 구간 180\,step이다.}
\label{fig:grpocredit}
\end{figure*}

\subsubsection{반사실적 크레딧 배분}

식~\eqref{eq:lpg}는 팀 보상의 이득 $A_{n,k}$ 를 모든 방어정의 로그확률에 같은 값으로
곱한다. 팀 보상을 모든 방어정이 그대로 나누어 가지면 개별 기여가 가려진다. 잘한 방어정과 못한
방어정이 같은 신호를 받으면 학습이 진행되지 않는다. 따라서 한 척의 행동만 바꾸고 나머지는
고정한 채 재평가하여, 그 차이를 해당 방어정의 몫으로 돌린다. 기준 표본
$\mathbf{a}^{\mathrm{base}} \sim \pi_\theta(\cdot \mid \mathbf{s})$ 를 하나 두고, 방어정 $p$ 에
대해 $p$ 의 행동만 바꾼 혼합 후보 $G$개를 만든다.
\begin{equation}
\begin{aligned}
\mathbf{a}^{(k, p)} &= \big(\mathbf{a}^{\mathrm{base}}_{-p},\; \tilde{a}^{(k)}_p\big),\\
\tilde{a}^{(0)}_p &= a_{\mathrm{heur},p}(\mathbf{s}), \qquad
\tilde{a}^{(k)}_p \sim \pi_\theta(\cdot \mid \mathbf{s}, p)
\end{aligned}
\label{eq:cfmix}
\end{equation}
여기서 $\mathbf{a}_{-p}$ 는 $p$ 를 제외한 나머지 방어정의 행동이다. 각 혼합 후보를 같은 구간
$H$ 만큼 재진행하여 얻은 창 보상 $R^{(p)}_k = R(\mathbf{s}, \mathbf{a}^{(k,p)}; H)$ 로 방어정
$p$ 만의 이득을 산정한다.
\begin{equation}
A^{(p)}_k = \frac{R^{(p)}_k - R^{(p)}_0}{\mathrm{std}(\{R^{(p)}_j\}_{j=0}^{G-1}) + \epsilon}
\label{eq:cfadv}
\end{equation}
나머지 방어정의 행동이 고정되어 있으므로 $R^{(p)}_k - R^{(p)}_0$ 는 방어정 $p$ 의 선택이
바뀐 데서만 온 차이이다. 이 이득은 방어정 $p$ 의 로그확률에만 곱해진다.
\begin{equation}
\mathcal{L}_{\mathrm{CF}}(\theta) = -\frac{\sum_{p} \sum_{n} w^{(p)}_n \sum_{k} A^{(p)}_{n,k}\,
  \log \pi_\theta\big(\tilde{a}^{(k)}_{n,p} \mid \mathbf{s}_n\big)}{G \sum_{p} \sum_n w^{(p)}_n}
\label{eq:lcf}
\end{equation}
$w^{(p)}_n$ 은 유효 마스크 $w_n$ 을 방어정 $p$ 의 후보 집합에 적용한 것이며, 이 모드에서
식~\eqref{eq:loss}의 엔트로피는 기준 표본 $\mathbf{a}^{\mathrm{base}}$ 에서만 계산한다.
식~\eqref{eq:lcf}는 식~\eqref{eq:loss}에서 $\mathcal{L}_{\mathrm{PG}}$ 자리에 들어간다. 이는
COMA~\citep{foerster2018coma}의 반사실적 기준선과 같은 착상이되, 학습된 비평가로 주변화하는
대신 \emph{시뮬레이터를 실제로 재진행}하여 얻는다는 점이 다르다. 비평가 근사 오차가 개입하지
않는 대신, 결정당 시뮬레이터 호출이 $G$ 회에서 $G P$ 회로 늘어난다.

\textbf{채택 여부.} 두 손실을 같은 예산(800 업데이트)으로 비교한 스크리닝에서 반사실
모드는 600 업데이트 시점 greedy 포획률 $0.801$ 로 팀 이득 모드의 $0.800$ 과 동률이었고,
시뮬레이터 호출은 3배였다. 본 논문이 평가하는 체크포인트는 \emph{팀 이득
손실}(식~\eqref{eq:lpg}--\eqref{eq:update})로 학습한 것이다.

\subsubsection{보상 설계}
\label{sec:reward}

한 후보의 창 보상 $R(\mathbf{s}, \mathbf{a}; H)$ 는 결정 시점부터 $H$\,step 동안 일어난
일을 하나의 스칼라로 합산한 팀 공유 보상이며, 이벤트 항·퍼텐셜 차분·조밀 항으로 구성된다.
\begin{equation}
R = R_\mathrm{event} + \big[\gamma\,\Phi(\mathbf{s}_{t+H}) - \Phi(\mathbf{s}_t)\big] + R_\mathrm{dense}
\label{eq:reward}
\end{equation}
이벤트 항은 창 안의 사건 수에 고정 가중을 곱한 합이다.
\begin{equation}
\begin{aligned}
R_\mathrm{event} = \;& r_{\mathrm{cap}} N_{\mathrm{cap}} + r_{\mathrm{brc}} N_{\mathrm{brc}}
  + r_{\mathrm{col}} \big(N_{\mathrm{col}} + N_{\mathrm{land}}\big)\\
  & + r_{\mathrm{tng}} N_{\mathrm{tng}}
  + \mathbb{1}_{\mathrm{end}}\big(r_{\mathrm{wipe}}\,\mathbb{1}[\mathcal{A} = \varnothing] + r_{\mathrm{surv}}\,|\mathcal{A}|\big)
\end{aligned}
\label{eq:revent}
\end{equation}
계수는 포획 $r_{\mathrm{cap}} = +13.5$, 돌파 $r_{\mathrm{brc}} = -8.0$, 아군·모선·육지 충돌
$r_{\mathrm{col}} = -30$, 설치 그물 접촉 $r_{\mathrm{tng}} = -13$, 창 안에서 에피소드가 끝날
때($\mathbb{1}_{\mathrm{end}}$) 전멸 $r_{\mathrm{wipe}} = +8$, 잔존 적 1척당 $r_{\mathrm{surv}} = -5$ 이다. $\Phi$ 는 적의
모선 근접도에 대한 퍼텐셜로, 차분 형태로 주어지므로 최적 정책을
보존한다~\citep{ng1999shaping}.
\begin{equation}
\Phi(\mathbf{s}) = -w_\mathrm{threat} \sum_{m \in \mathcal{A}}
          \mathrm{clip}\!\left(1 - \frac{\lVert \mathbf{e}_m - \mathbf{m} \rVert}{R_w},\, 0,\, 1\right)
\label{eq:potential}
\end{equation}
여기서 $\mathcal{A}$ 는 생존 적선 집합, $R_w = 6300\unit{m}$ 는 교전 해역 반폭,
$w_\mathrm{threat}=3.0$, $\gamma = 0.99$ 이다. 적이 포획되어 $\mathcal{A}$ 에서 빠지거나
멀어지면 $\Phi$ 가 오른다.

조밀 항은 창 안의 기동 품질을 재는 비용·보너스의 합이다. 창 길이 $H$, 방어정 속력 $v_a$,
최대 선회율 $\dot\psi_{\max}$, 창 안 총 이동거리 $D$, 소모 그물 수 $N_{\mathrm{net}}$ 에 대해
\begin{equation}
\begin{aligned}
R_\mathrm{dense} = \;& -\,w_{\mathrm{path}} \frac{D}{H v_a P} - w_{\mathrm{net}} N_{\mathrm{net}}
  - \frac{w_{\mathrm{turn}}}{P}\sum_{p=1}^{P} S_p\\
  & - w_{\mathrm{prox}} \sum_{p=1}^{P}\big(\rho_{p}^{\mathrm{m}\,2} + \rho_{p}^{\mathrm{a}\,2}\big)
  - \frac{w_{\mathrm{cons}}}{P}\sum_{p=1}^{P} \Xi_p\\
  & + w_{\mathrm{cov}}\, \mathrm{Cov} + w_{\mathrm{ray}}\, \mathrm{Ray}
\end{aligned}
\label{eq:rdense}
\end{equation}
이다. $S_p = \sum_t |\Delta\psi_{p,t}| / (H \dot\psi_{\max}) \in [0,1]$ 은 창 안 누적 선회량을
최대 가능 선회량으로 나눈 값이고, $\Xi_p = \lVert \mathbf{t}_p - \mathbf{t}^{\mathrm{prev}}_p \rVert / L_{\max}$
는 직전 결정 대비 목표점 이동을 최대 경로 길이 $L_{\max} = L_{\mathrm{net}}$ 로 나눈 값이다(구현은
경유점 6개와 생존 방어정에 대해 평균하며, 여기서는 목표점 단위로 단순화하여 적었다). 근접 페널티는 창 안에서
방어정 $p$ 가 모선·타 방어정에 가장 가까웠던 거리
$d^{\mathrm{m}}_p$, $d^{\mathrm{a}}_p$ 를 영향 반경으로 정규화한 뒤 제곱하여, 붙을수록
급격히 커지게 한다.
\begin{equation}
\rho^{\mathrm{m}}_p = \mathrm{clip}\!\Big(\frac{R_{\mathrm{m}} - d^{\mathrm{m}}_p}{\varrho_{\mathrm{m}}}, 0, 1\Big), \quad
\rho^{\mathrm{a}}_p = \mathrm{clip}\!\Big(\frac{\varrho_{\mathrm{a}} - d^{\mathrm{a}}_p}{\varrho_{\mathrm{a}}}, 0, 1\Big)
\label{eq:prox}
\end{equation}
($R_{\mathrm{m}} = 800$, $\varrho_{\mathrm{m}} = 500$, $\varrho_{\mathrm{a}} = 600\unit{m}$).
커버리지는 적에서 모선으로 향하는 직선(ray)이 설치된 그물에 막힌 생존 적의 비율이며,
포획의 선행 지표로서 창 안에 포획이 없을 때도 변별력을 준다.
\begin{equation}
\mathrm{Cov} = \frac{1}{|\mathcal{A}|} \sum_{m \in \mathcal{A}}
  \mathbb{1}\big[\exists\, \lambda \in \Lambda :\; \mathbf{e}_m + \lambda(\mathbf{m} - \mathbf{e}_m) \in \mathcal{W}\big]
\label{eq:coverage}
\end{equation}
여기서 $\mathcal{W}$ 는 설치된 그물이 차지하는 영역이고, $\Lambda$ 는 $[0.05, 0.95]$ 를
16등분한 표본점이다. $\mathrm{Ray}$ 는 같은 검사를 계획된 그물벽까지 포함하되, 두 방어정이 충돌 거리 안에서
같은 ray 를 겹쳐 막으면 무효로 치고, 차단된 적마다 배율
$b_{\mathrm{net}}\,(1 + g_{\mathrm{far}}(1 - \lambda_m))$ ($b_{\mathrm{net}} = 1.3$,
$g_{\mathrm{far}} = 1.0$, $\lambda_m$ 은 차단점의 ray 상 위치로 적에 가까울수록 0)을 곱해
적에 가까운 차단을 우대하는 단일 커버 판이다($\mathrm{Ray} \in [0, 2.6]$). 가중은 $w_{\mathrm{path}} = 0.5$,
$w_{\mathrm{net}} = 0.3$, $w_{\mathrm{prox}} = 0.75$, $w_{\mathrm{turn}} = w_{\mathrm{cons}} = 0.3$,
$w_{\mathrm{cov}} = 0.1$, $w_{\mathrm{ray}} = 1.6$ 이다. 레이 커버 항을 제외한 조밀 항은 이벤트
계수의 수십분의 일 규모이고, 레이 커버 항은 최대 $4.2$ 로 포획 1건($13.5$)의 3분의 1을 넘지
않게 두었다. 이 밖에 요격점 접근·예비 방어정 정지·인공 퍼텐셜
척력 등 가중 $0.1\sim0.25$ 의 보조 항이 있으나 규모가 작아 생략한다.

퍼텐셜 항을 둔 이유는 이벤트 항만으로는 학습 신호가 지나치게 희소하기 때문이다. 216개
결정에 대한 실측에서 포획은 평가 창의 46.8\,\%에서만 발생하였다. 나머지 절반의 창에서는
퍼텐셜 항이 유일한 큰 신호이며(비영 비율 100\,\%), 이것이 없으면 그룹 내 후보들의 순위가
갈리지 않아 식~\eqref{eq:adv}의 이득이 정의되지 않는다. 같은 측정에서 퍼텐셜 항의 평균은
$+0.36$ 인 데 비해 절댓값 평균은 $7.40$ 으로, 창마다 크게 변동하지만 전체적으로는 상쇄된다.
이는 퍼텐셜 기반 성형이 설계대로 \emph{정책은 바꾸지 않고 변별력만 제공}한다는 실측 증거이다.

\subsubsection{수렴}

학습은 실해역 6곳을 에피소드마다 돌려가며 진행하였다(Fig.~\ref{fig:sites}). 한 해역만으로
학습하면 정책이 그 해안선의 형상을 외워 ``섬은 항상 저기 있다''를 전제로 굳는다. 여섯 곳은
육지 비율 0\,\%(외해)에서 7.5\,\%까지, 요격 환형에 육지가 잠식하는 비율 0에서 6\,\%까지
분포하도록 골랐고, 모선에서 가장 가까운 육지까지의 여유도 1.6\,km에서 무한대까지 걸쳐 있다.
평가(\S\ref{sec:experiment})는 이 가운데 통영 매물도 서측 한 곳에서 수행한다.

%-------------------------------------------------- 학습 해역 6곳
\begin{figure*}[t]
\centering
\includegraphics[width=0.90\textwidth]{fig_sites6.png}
\caption{학습에 쓴 실해역 6곳의 육지 마스크(Fig.~\ref{fig:landpipe}의 절차를 각 해역에
         적용한 결과, 12.6\,km 박스). 별표는 모선, 파선 원은 요격 환형이다. 각 패널 위의
         백분율은 박스 안 육지 비율, 아래는 요격 환형에 육지가 잠식한 비율과 모선에서 가장
         가까운 육지까지의 여유 거리이다. 오른쪽 아래(포항 호미곶 동측)는 육지가 없는 외해로,
         지형이 없는 상황을 학습 분포에 넣기 위한 대조 해역이다.}
\label{fig:sites}
\end{figure*}

같은 여섯 해역에서 학습된 정책이 실제로 어떻게 움직이는지는 Fig.~\ref{fig:sitesrun}에
같은 진행 시점(에피소드의 약 65\,\%, 그물이 전개되고 적이 요격 환형에 도달한 시점)으로
나란히 놓았다. 집중 포메이션·휴리스틱 배정·U-Net 기동으로 통일하였고 해역만 다르다. (c)
신안 홍도 북측과 (e) 통영 매물도 북측에서는 해안선이 요격 환형 안쪽까지 들어와 유효
마스크를 잠식하며, (f) 포항 호미곶 동측은 육지가 없어 정책이 지형 제약 없이 기동하는 대조
사례다. 각 패널의 시드는 해당 해역에서 전멸 포획·돌파 0·충돌 0·아군 손실 0을 만족하는
것으로 골랐다 --- 이 그림은 성능 주장이 아니라 지형이 경로와 그물 위치에 어떻게 개입하는지를
보이는 용도이다.

\begin{figure*}[t]
\centering
\includegraphics[width=\textwidth]{fig_sites6_run.png}
\caption{학습 해역 6곳의 실행 장면(집중 포메이션, 휴리스틱 배정, U-Net 기동, 에피소드
         진행률 65\,\%). 패널 순서와 육지 비율은 Fig.~\ref{fig:sites}와 같다. 위성 배경 위에
         육지 마스크·설치된 그물(보라)·경유점·경로가 겹쳐 있고, 각 패널 상단에 전장
         상황판과 적/아군 확대 창이 있다. 시드는 해역별로 사고 없는 에피소드를 골랐다
         (\texttt{tools/find\_demo\_seeds.py --site}).}
\label{fig:sitesrun}
\end{figure*}

Fig.~\ref{fig:traincurve}에 학습 경과를 정리하였다. 포획률(패널 (a))은 휴리스틱 기준선
0.744 부근에서 출발하여 800 업데이트 안에 greedy 평가 최고 0.805(600 업데이트)에 도달한다.
본 논문이 평가하는 정책은 \textbf{이 최고 시점의 가중치}이다.

세 가지를 함께 볼 필요가 있다. 첫째 \textbf{보상(패널 (b))에서 후보 최대값이 후보 평균보다
꾸준히 위에 있다} --- 그룹 안에서 후보들이 실제로 갈린다는 뜻이며, 갈리지 않으면
식~\eqref{eq:adv}의 이득이 0이 되어 학습 신호가 사라진다. 둘째 \textbf{학습신호 유효율
(패널 (c))이 0.5 안팎에 머문다} --- 절반가량의 월드는 후보 간 분산이 0이어서 그 스텝의
기울기에 기여하지 못한다. 셋째 손실(패널 (d))의 절대값은 정책경사와 엔트로피 항의 합이므로
그 자체로는 수렴 지표가 아니다.

포획률 곡선이 단조 증가하지 않고 진동하는 것은 이 구조의 특성이다. 매 업데이트마다 새로
표집한 전장에서 평가하므로 시나리오 난이도가 곡선에 그대로 실린다. 그래서 학습 중 판정은
greedy 평가(패널 (a)의 원 표식)로 하고, 최종 성능은 \S\ref{sec:experiment}의 고정 시드
대응표본으로 따로 측정한다.

\begin{figure*}[t]
\centering
\includegraphics[width=0.92\textwidth]{fig_traincurve.pdf}
\caption{GRPO 학습 경과($50^2$ 격자, 후보 $G=8$, 후보 평가 구간 180\,step, 실해역 6곳
         로테이션(Fig.~\ref{fig:sites}), 800 업데이트). (a) 포획률. 옅은 선은 학습 중 후보 평균, 원 표식은
         200 업데이트마다의 greedy 평가, 회색 파선은 휴리스틱 기준선 0.744이다.
         \textbf{별표(600 업데이트, 0.805)가 본 논문이 평가하는 가중치}이다.
         (b) 후보 평균 보상과 후보 최대 보상 --- 둘의 간격이 그룹 내 변별력이다.
         (c) 후보 간 분산이 0이 아닌 월드의 비율(학습 신호가 실제로 생긴 비율).
         (d) 손실(정책경사 $+$ 엔트로피). 부호와 크기 자체는 수렴 지표가 아니다.}
\label{fig:traincurve}
\end{figure*}

\paragraph{수렴 진단}
\label{sec:convergence}
Fig.~\ref{fig:traincurve}는 채택 런의 기록이지만, 그 런은 보상·포획률·손실만 20 업데이트마다
남겼으므로 \emph{정책이 정지점에 닿았는가}를 직접 보이는 지표가 없다. 이를 보완하기 위해
같은 설정을 시드 3개로 다시 돌리며 표준적인 수렴 진단을 매 업데이트마다 기록하였다
(Fig.~\ref{fig:convergence}). 재현 런의 가중치는 논문의 결과에 쓰지 않는다 --- 채택 런이
\emph{전형적인} 학습 경과였는지, 그리고 그 경과가 수렴이라 부를 만한지를 확인하는 용도이다.
갱신 전후 정책을 $\theta_{t}$, $\theta_{t+1}$ 라 하면 두 진단량은 다음과 같다.
\begin{equation}
\begin{aligned}
\widehat{\mathrm{KL}}_t &= \frac{1}{NP(G-1)}\sum_{n,p}\sum_{k=1}^{G-1}
  \big(\zeta_{n,p,k} - 1 - \log \zeta_{n,p,k}\big),\\
\zeta_{n,p,k} &= \frac{\pi_{\theta_{t+1}}\big(a^{(k)}_{n,p} \mid \mathbf{s}_n\big)}{\pi_{\theta_t}\big(a^{(k)}_{n,p} \mid \mathbf{s}_n\big)}, \qquad
\varsigma_t = \frac{\lVert \theta_t - \theta_{\mathrm{BC}} \rVert}{\lVert \theta_{\mathrm{BC}} \rVert}
\end{aligned}
\label{eq:diag}
\end{equation}
$\widehat{\mathrm{KL}}_t$ 는 같은 관측·같은 표본 후보(휴리스틱 후보 $k=0$ 은 정책의 표본이
아니므로 제외)에 대한 확률비 $\zeta$ 의 $k_3$ 추정량으로,
$\mathrm{KL}(\pi_{\theta_t} \,\|\, \pi_{\theta_{t+1}})$ 의 비음(非陰)·저분산 몬테카를로 추정이다.
$\varsigma_t$ 는 행동복제 초기점 $\theta_{\mathrm{BC}}$ 로부터의 상대 파라미터 이동 거리이다.

네 지표가 같은 방향을 가리킨다. (i) \textbf{업데이트 간 근사 KL}(패널 (c)) --- 식~\eqref{eq:diag}의
$\widehat{\mathrm{KL}}_t$ --- 은 처음 100 업데이트 평균
$1.4\times10^{-2}$ 에서 마지막 100 업데이트 $1.5\times10^{-4}$ 로 두 자릿수 줄어든다. 같은
구간에서 학습률은 16배 줄었으므로 이 감소의 일부는 스케줄의 몫이지만, 정책이 한 걸음에
움직이는 양이 사실상 0 에 수렴한다는 사실 자체는 변하지 않는다. (ii) \textbf{파라미터 이동
거리} $\varsigma_t$(패널 (f))는 400
업데이트에 $0.031$, 800 업데이트에 $0.034$ 로 뒤 절반에서 포화한다 --- 정책이 행동복제
초기점 근방의 한 지점에 정착했다는 뜻이다. (iii) \textbf{정책 엔트로피}(패널 (b); 1단계 픽셀
분포, 균일분포 상한 $\ln 273 \approx 5.6$\,nat)는 $2.42$ 에서 $2.06$\,nat 로 내려가 400
업데이트 이후 평탄하고, 최빈 픽셀의 확률은 $0.25$ 에서 $0.37$ 로 오른다. 정책이 붕괴하지
않으면서(엔트로피가 0 으로 가지 않으면서) 첨예해진다. (iv) \textbf{휴리스틱 후보 대비 보상
이득}(패널 (e)) --- 같은 전장·같은 평가 창에서 표본 후보의 보상에서 휴리스틱 후보의 보상을 뺀
값 --- 은 처음 100 업데이트 평균 $-1.3$ 에서 마지막 100 업데이트 $+5.1$ 로 부호가 바뀐다.
학습 초기에는 정책의 표본이 휴리스틱보다 못했고, 후반에는 평균적으로 휴리스틱을 앞선다.

greedy 평가 포획률(패널 (a))은 세 시드 모두 800 업데이트에서 기준선 $0.749$ 위에 있다
($0.817$, $0.777$, $0.801$; 평균 $0.798$). 채택 런의 세 평가점(별표)은 이 세 곡선의 범위
안에 들어, 채택 런이 특별히 운 좋은 경과가 아니었음을 보인다. 다만 시드 하나는 100
업데이트에서 $0.417$ 까지 떨어졌다가 회복하였고 500 업데이트에서 다시 기준선 아래로
내려갔다 --- 그룹 상대 이득이 순위만 쓰는 이 학습이 초기에 한 번 크게 흔들릴 수 있음을 그대로
보여 주는 사례이며, 학습 중 판정을 greedy 평가로 하는 이유이기도 하다.

한 가지 지표는 수렴 지표로 읽으면 안 된다. \textbf{그래디언트 노름}(패널 (d))은 800
업데이트 내내 클립 한계 1.0 의 12배 이상이어서 모든 갱신이 잘린 상태로 진행된다. 이는
식~\eqref{eq:adv}의 이득이 그룹 표준편차로 정규화되어 크기가 없는 양이기 때문이다 --- 후보들의
보상 차이가 아무리 작아도 이득의 스케일은 같으므로 기울기의 크기는 줄지 않는다. 따라서 이
구조에서 실제 보폭을 정하는 것은 기울기가 아니라 학습률이며, 수렴의 증거는 기울기 노름이
아니라 KL·파라미터 이동·엔트로피에서 찾아야 한다.

\begin{figure*}[t]
\centering
\includegraphics[width=0.96\textwidth]{fig_convergence.pdf}
\caption{수렴 진단 --- 채택 런과 같은 설정을 시드 3개로 다시 돌린 학습 동역학. 굵은 선은
         시드 평균, 띠는 시드 간 최소~최대, 시계열은 20 업데이트 이동평균이다.
         (a) 100 업데이트마다의 greedy 포획률(옅은 선은 시드별, 원은 평균, 별표는 채택 런의
         평가점, 회색 파선은 같은 프로토콜로 잰 휴리스틱 기준선 0.749).
         (b) 1단계 픽셀 분포의 엔트로피. 파선은 유효 픽셀 중앙값 273개의 균일분포.
         (c) 갱신 전후 정책의 근사 KL 발산(같은 관측·같은 표본 후보, 로그축).
         (d) 클리핑 전 그래디언트 노름(로그축) --- 항상 클립 한계 위에 있어 수렴 지표가 아니다.
         (e) 표본 후보 보상에서 휴리스틱 후보 보상을 뺀 값. 0 을 넘으면 정책의 표본이
         휴리스틱을 앞선다.
         (f) 행동복제 초기점으로부터의 상대 파라미터 이동 거리.}
\label{fig:convergence}
\end{figure*}

\subsubsection{초기 학습 안정화}
\label{sec:bc}

그룹 상대 이득은 후보들이 실제로 갈릴 때만 신호를 주므로, 무작위 초기화된 정책에서 곧바로
시작하면 초기 후보가 모두 무의미한 자리에 흩어져 갱신이 잡히지 않는다. 그래서 그룹 상대
최적화에 앞서 휴리스틱 행동(\S\ref{sec:heuristic})을 목표로 하는 행동 복제를 300 step
수행하여 정책을 부팅한다(Adam, 학습률 $3\times10^{-4}$). 목표 분포에는 가우시안 라벨
스무딩을 적용한다. 휴리스틱이 지목한 픽셀 $i^{*}_k$ 주변에 유효 픽셀로 제한한 가우시안
가중을 두고, 그 가중과 정책의 교차엔트로피를 최소화한다.
\begin{equation}
\omega_k(i) = \frac{\mathbb{1}[i \in \mathcal{V}^{(k)}_p]\,
  \exp\!\big(-\lVert \mathbf{x}_i - \mathbf{x}_{i^{*}_k} \rVert^2 / 2\sigma_{\mathrm{BC}}^2\big)}
  {\sum_{j \in \mathcal{V}^{(k)}_p} \exp\!\big(-\lVert \mathbf{x}_j - \mathbf{x}_{i^{*}_k} \rVert^2 / 2\sigma_{\mathrm{BC}}^2\big)}
\label{eq:bcweight}
\end{equation}
\begin{equation}
\mathcal{L}_{\mathrm{BC}}(\theta) = -\frac{1}{P}\sum_{p=1}^{P} \sum_{k=0}^{K-1} \sum_{i \in \mathcal{V}^{(k)}_p}
  \omega_k(i)\, \log \pi^{(k)}_p(i)
\label{eq:lbc}
\end{equation}
$\sigma_{\mathrm{BC}} = 0.8$\,px($\approx 202\unit{m}$)이며, 두 번째 픽셀의 질의는 휴리스틱의
첫 픽셀로 갱신한다(teacher forcing). 격자 위의 인접 픽셀은 사실상 같은 행동이므로 단일
픽셀 원-핫을 목표로 두면 불필요하게 뾰족한 분포를 강요하게 되고, 조밀한 점수맵에서
원-핫 교차엔트로피의 기울기는 지나치게 희소하여 초기 학습이 붙지 않기 때문이다. 그룹 상대
최적화 단계에서는 이 항을 쓰지 않으며, 대신 식~\eqref{eq:adv}의 휴리스틱 기준선이 정책을
휴리스틱 근방에 붙잡는 신뢰 영역 구실을 한다.

\subsubsection{학습률과 엔트로피 계수의 스케줄}
\label{sec:schedule}

옵티마이저는 Adam이고 학습률은 재시작 없는 단일 주기 코사인
어닐링~\citep{loshchilov2017sgdr}으로 줄인다(Fig.~\ref{fig:lrsched}(a)):
\begin{equation}
\begin{aligned}
\eta_t &= \eta_{\min} + \tfrac{1}{2}\,(\eta_0 - \eta_{\min})
          \Bigl(1 + \cos\tfrac{\pi t}{T}\Bigr),\\
\eta_0 &= 10^{-4},\quad \eta_{\min} = 0.05\,\eta_0,\quad T = 800 .
\end{aligned}
\label{eq:cosine}
\end{equation}
주기 $T$를 총 업데이트 수에 묶었으므로 예산을 바꾸면 스케줄 전체가 함께 늘어난다 --- 비교
실험은 같은 $T$로만 한다. 같은 축 위에서 식~\eqref{eq:loss}의 엔트로피 계수는 선형으로
감쇠한다(Fig.~\ref{fig:lrsched}(b)).
\begin{equation}
\lambda_{\mathrm{H}}(t) = \lambda_0 \Big(1 - (1 - \kappa)\frac{t}{T}\Big), \qquad
\lambda_0 = 0.01,\quad \kappa = 0.1
\label{eq:entsched}
\end{equation}
즉 $0.01$ 에서 $0.001$ 로 줄어든다. 기울기는 식~\eqref{eq:update}대로 노름 $g_{\max} = 1.0$ 에서
자른다.

코사인을 고른 이유는 이 학습 구조의 두 가지 성질 때문이다. 첫째, 그룹 상대 이득
(식~\eqref{eq:adv})은 매 업데이트마다 새로 표집한 전장에서 $G=8$개 후보의 \emph{순위}만
쓰므로 기울기의 분산이 크고 시나리오 난이도가 그대로 실린다(Fig.~\ref{fig:traincurve}(a)의
진동). 앞 구간의 큰 학습률은 휴리스틱 부팅 분포에서 빨리 벗어나게 하고, 뒤 구간의 작은
학습률은 이 잡음 섞인 기울기가 정책을 이리저리 끌지 못하게 한다. 둘째, 학습률과 엔트로피
계수를 \emph{같은 예산 축}에서 함께 줄이면 분포가 뾰족해지는 시점과 보폭이 작아지는 시점이
맞물려, 탐색이 남아 있는 동안에는 큰 걸음을, 수렴 국면에는 작은 걸음을 딛는다. 계단식
감쇠와 달리 불연속점이 없어 정책의 표본 분포가 갑자기 바뀌지 않는 점도 순위 기반 이득에는
유리하다.

학습 경과에서 이 설계와 일치하는 관찰은 다음과 같다. 학습 중 후보 평균 포획률을 세 구간으로
나누면 앞 두 구간(업데이트 20--260, 280--520)은 평균 $0.697$, $0.703$ 에 표준편차 $0.060$,
$0.069$ 인 데 비해, 학습률이 $\eta_0$의 20\,\% 아래로 내려간 마지막 구간(540--800)은 평균
$0.762$, 표준편차 $0.048$ 로 더 높고 더 좁다. 채택한 가중치(greedy 평가 최고 0.805)도 이
구간의 600 업데이트 시점, $\eta_t \approx 1.9\times10^{-5}$ 에서 나왔다. 다만 상수 학습률과의
대조 실험은 수행하지 않았으므로, 이 관찰은 스케줄의 효과를 분리해 보인 것이 아니라 설계
의도와 모순되지 않음을 보인 것이다.

\begin{figure}[t]
\centering
\includegraphics[width=\linewidth]{fig_lr_schedule.pdf}
\caption{학습 코드가 실제로 쓴 스케줄. (a) 식~\eqref{eq:cosine}의 단일 주기 코사인 학습률.
         원 표식은 200 업데이트마다의 greedy 평가 시점, 별표는 채택 가중치(600 업데이트,
         $\eta_t \approx 1.9\times10^{-5}$), 회색 파선은 $\eta_{\min}$이다.
         (b) 엔트로피 계수 $\lambda_{\mathrm{H}}(t)$의 선형 감쇠($0.01 \to 0.001$).}
\label{fig:lrsched}
\end{figure}
\subsection{전략 계층: 조합 최적화 없는 LLM 지휘관}
\label{sec:commander}

전략 계층의 위치는 Fig.~\ref{fig:architecture}에, 실제 동작 화면은
Fig.~\ref{fig:runshot}에 있다. 입력은 전장 상태·임무 목표·교전 규칙·자산 상태이고,
출력은 무리별 배정과 대기(HOLD)이며, 그 배정이 기동 계층으로 전달되어 각 방어정의
그물벽 위치 결정을 조건짓는다.

언어모델에는 좌표 원시 데이터 대신 \emph{기하량이 미리 계산된} 전장 상태를 직렬화하여
전달한다. 각 적 무리의 요격점까지의 이동거리와 선회량, 모선 기준
방위, 설치된 그물에 의한 접근로 차단 여부 등이 포함된다. 언어모델이 삼각함수를 암산하게
하는 대신 판단에 필요한 양을 미리 제공하는 것이다.

\subsubsection{구조적 출력에 의한 행동 규제}
\label{sec:schema}

언어모델의 출력을 자유 문장으로 받으면 하위 계층이 파싱에 실패할 수 있고, 실패 시 그
주기의 배정이 통째로 사라진다. 본 연구는 출력을 \emph{스키마로 강제}하여 이 실패 모드를
구조적으로 차단한다~\citep{openai2026function}. 계획을 다음 자료형으로 선언하고, 그로부터 생성한
JSON 스키마를 모델 호출에 함께 넘긴다.

% 좁은 단에 맞춰 줄을 미리 끊는다. 자동 줄바꿈에 맡기면 들여쓰기가 무너져
% 자료형 구조가 안 보인다.
\begin{quote}\ttfamily\scriptsize\setlength{\parindent}{0pt}
CommanderPlan \{\\
\hspace*{1.2em} rationale: str\\
\hspace*{1.2em} deployments: List[\{\\
\hspace*{3.0em} cluster\_id: int,\\
\hspace*{3.0em} ally\_ids: List[int] \}]\\
\hspace*{1.2em} hold\_ships: List[int]\\
\}
\end{quote}

이 스키마는 디코딩 단계에서 문법으로 작용하여, 형식에 맞지 않는 토큰이 애초에 생성되지
않게 한다. 지역 실행에서는 스키마를 \texttt{format} 인자로, 상용 API에서는
\texttt{response\_schema} 로 전달하며, 두 경우 모두 응답을 같은 자료형으로 역직렬화하여
검증한다. 온도는 0으로 두어 같은 전장에서 같은 배정이 나오게 한다. 이는 \S\ref{sec:paired}의
시드 대응표본 설계가 성립하기 위한 전제이기도 하다.

\textbf{필드 순서가 곧 사고 순서다.} 스키마는 \emph{판단 근거 $\to$ 무리별 배정 $\to$ 대기
목록} 순으로 고정된다. 자기회귀 생성에서 뒤 필드는 앞 필드를 조건으로 하므로, 근거를 먼저
쓰게 하면 배정이 그 근거를 조건으로 결정된다. 순서를 뒤집어 배정을 먼저 내게 하면 근거는
이미 정해진 답에 대한 사후 설명이 되므로, 최소한 설명 가능성은 잃는다. 실제로 이
필드에는 각 방어정의 요격점까지 거리와 선회량을 인용하도록 요구하여, 근거가 제공된 기하량에
닿아 있게 만들었다. \textbf{다만 이 요구의 준수 여부는 본 연구가 검증하지 못한다} --- 실험에서
보존한 것은 근거의 \emph{길이}뿐이고 원문은 저장하지 않았기 때문이다(\S\ref{sec:limitations}).
따라서 이 설계는 근거를 확인 가능하게 \emph{만들 수 있는} 구조라는 주장이지, 본 실험이 그
확인을 수행했다는 주장이 아니다.

\textbf{빈 배정의 의미는 프롬프트가 규정한다.} \texttt{ally\_ids} 를 비우는 것은 ``시스템이
대신 고르라''가 아니라 ``이 무리에는 아무도 보내지 않는다''를 뜻한다. 앞에서 밝혔듯 코드
보완이 없으므로 빈 목록은 그대로 무방비를 낳으며, 프롬프트가 이 결과를 명시하여
지휘관이 그 대가를 알고 결정하게 한다.

\textbf{조합 최적화 후처리를 두지 않는다.} 언어모델이 지목한 배가 그대로 배정되며, 최소비용
매칭이나 2-opt 개선이 개입하지 않는다. 언어모델이 비워 둔 무리는 담당 없이 남고, 지목되지
않은 방어정은 대기한다. 후처리를 두면 최종 배정이 사실상 비용 함수의 산물이 되어 언어모델의 기여를 측정할 수 없기 때문이다. 이 선택은 성능이 아니라
\emph{측정 가능성}을 위한 것이며, 그 대가는 \S\ref{sec:results}에서 정량적으로 드러난다.

\textbf{후처리를 두지 않는다는 것이 기하 보정이 없다는 뜻은 아니다.} 언어모델이 정할 수
있는 것은 출력 스키마가 정확히 제한한다 --- 무리별 배정 항목의 필드는 \texttt{cluster\_id} 와
\texttt{ally\_ids} 둘뿐이고, 그 밖에 결정되는 것은 판단 근거와 대기 목록(\texttt{hold\_ships})
뿐이다. 즉 권한은 \emph{어느 배가 어느 무리를 맡는가}와 \emph{이번 주기에 누가 멈추는가}로
닫혀 있다. 그물 전개 여부도, 벽의 방위도, 교전 반경도 출력에 없다. 배정으로부터 요격점을
만드는 기하는 전적으로 코드다.

그 기하는 두 경로에서 같은 규칙을 쓴다. 한 무리에 두 척 이상이 배정되면
\S\ref{sec:heuristic}의 잉여 배정과 \emph{같은 두 손잡이}가 작동한다 --- 모선 쪽으로
$\lambda_{\mathrm{L}} = 0.18$ 씩 미는 반경 층 분리와 $\pm 40^\circ$ 씩 번갈아 트는 방위 분리다. 완전히
같지는 않아서, 언어모델 경로는 (i) 층 순서를 기본 요격점에 가까운 배부터 주고, (ii) 층을
깊게 넣다 요격점이 모선 코앞으로 들어오는 것을 막기 위해 반경을 요격 환형 안으로 자른다.
환형을 벗어난 요격점은 기동 계층의 유효 마스크를 붕괴시켜 폴백으로 떨어지기 때문이다.

권한이 배정과 대기로 닫혀 있다는 점이 중요하다. 전략 계층이 기하까지 건드릴 수 있으면
\S\ref{sec:res_maneuver}에서 관측되는 차이가 배정 판단에서 온 것인지 기하 조정에서 온 것인지
분리되지 않는다. 두 경로가 같은 기하 규칙을 공유하기 때문에, 조건 간 차이는 배정 그 자체의
차이로 읽을 수 있다.

배정이 확정되면 각 방어정의 요격점이 정해지고, 그 요격점을 중심으로 기동 계층이 그물벽을 놓을
수 있는 후보 영역(유효 마스크)이 함께 결정된다. 즉 배정은 기동 계층의 선행 조건이며,
배정이 바뀌면 후보 영역 자체가 바뀐다.
\subsection{배치 관점}
\label{sec:deploy}

제안 정책에는 배치 통계에 의존하는 정규화가 일절 없다. 정규화는 전부 그룹 정규화이고
(\S\ref{sec:policy}), 관측은 격자 크기와 설정 상수로만 나누어 정규화한다 --- 이동 평균으로
갱신되는 running statistics 가 없다는 뜻이다. 그래서 배치 크기 1의 실시간 추론이 학습과 같은
연산을 수행한다. 통계 워밍업이 필요 없고, 학습 모드와 추론 모드의 동작이 갈리지 않는다.

또한 정책은 경유점만 지정하고 조타는 저수준 제어기가 담당하므로 선체 동역학과 분리되어
있다. 수상정의 조종 동역학은 확립된 모형이 존재하므로~\citep{fossen2011marine} 이 분리는
정책을 재학습하지 않고도 선체 특성이 다른 플랫폼에 이식할 여지를 남긴다. 관측과 행동은
절대 길이가 아닌 격자 상대 좌표로 표현된다. 이는 실해역 이식 시
격자 축척만 재설정하면 되는 구조로, Sim2Real 전이에 유리한 성질이다. 다만 본 논문은
시뮬레이션 검증에 한정되며, 이 성질의 실증은 \S\ref{sec:limitations}에서 논의한다.
